\documentclass{beamer}
\usetheme{AnnArbor}

\usepackage[T1]{fontenc}
\usepackage[utf8]{inputenc}
\usepackage[spanish]{babel}
\usepackage{amsmath}
\usepackage{tikz}
\usetikzlibrary{angles,quotes,calc}

\title{Animación}
\subtitle{Transformación: Animación}
\author{HCI}
\date{April 2026}

\begin{document}

%---------------------------------
\begin{frame}
\titlepage
\end{frame}

%---------------------------------
\begin{frame}{Objetivo}
Derivar matemáticamente la transformación de rotación:
\[
\begin{bmatrix}
x' \\
y'
\end{bmatrix}
=
\begin{bmatrix}
\cos\theta & -\sin\theta \\
\sin\theta & \cos\theta
\end{bmatrix}
\begin{bmatrix}
x \\
y
\end{bmatrix}
\]

\medskip

Y comprender formalmente por qué el empleo de las funciones seno y coseno.
\end{frame}

%---------------------------------
\begin{frame}{Punto en coordenadas cartesianas}
Un punto en el plano se representa como:

\[
P = (x,y)
\]

\begin{center}
\begin{tikzpicture}[scale=0.8]
\draw[->] (-0.5,0) -- (5,0) node[right] {$x$};
\draw[->] (0,-0.5) -- (0,3.5) node[above] {$y$};

\coordinate (P) at (4,2.5);

\draw[dashed] (P) -- (4,0);
\draw[dashed] (P) -- (0,2.5);

\filldraw[blue] (P) circle (2pt);
\node[above right] at (P) {$P=(x,y)$};
\end{tikzpicture}
\end{center}

Interpretación:
\begin{itemize}
\item \(x\): desplazamiento horizontal.
\item \(y\): desplazamiento vertical.
\end{itemize}
\end{frame}

%---------------------------------
\begin{frame}{Cambio a coordenadas polares}
El mismo punto puede describirse como:

\[
P = (r,\alpha)
\]

Donde:

\begin{itemize}
\item \(r\): distancia desde el origen hasta el punto.
\item \(\alpha\): ángulo respecto al eje \(x\).
\end{itemize}

\medskip

Este cambio es clave para entender la rotación.
\end{frame}

%---------------------------------
\begin{frame}{Interpretación geométrica}
\begin{center}
\begin{tikzpicture}[scale=0.8]
\draw[->] (-0.5,0) -- (5,0) node[right] {$x$};
\draw[->] (0,-0.5) -- (0,3.5) node[above] {$y$};

\coordinate (P) at (4,2.5);

\draw[thick,blue] (0,0) -- (P) node[midway,above] {$r$};

\draw[->,red,thick] (1,0) arc (0:32:1);
\node[red] at (1.3,0.3) {$\alpha$};

\filldraw[blue] (P) circle (2pt);
\node[above right] at (P) {$P$};
\end{tikzpicture}
\end{center}

El punto se localiza mediante una distancia \(r\) y un ángulo \(\alpha\).
\end{frame}

%---------------------------------
\begin{frame}{Triángulo rectángulo asociado}
\begin{center}
\begin{tikzpicture}[scale=0.85]
\coordinate (O) at (0,0);
\coordinate (X) at (4,0);
\coordinate (P) at (4,2.5);

\draw[->] (-0.5,0) -- (5,0) node[right] {$x$};
\draw[->] (0,-0.5) -- (0,3.2) node[above] {$y$};

\draw[thick] (O)--(X)--(P)--cycle;

\node[below] at (2,0) {$x$};
\node[right] at (4,1.25) {$y$};
\node[above left] at (2,1.25) {$r$};

\draw[->,red,thick] (0.9,0) arc (0:32:0.9);
\node[red] at (1.2,0.35) {$\alpha$};

\filldraw[blue] (P) circle (2pt);
\end{tikzpicture}
\end{center}

Se forma un triángulo rectángulo donde:

\[
r^2 = x^2 + y^2
\]
\end{frame}

%---------------------------------
\begin{frame}{Origen del coseno}
\small
El coseno se define como la razón entre el cateto adyacente y la hipotenusa:

\[
\cos\alpha = \frac{x}{r}
\quad \Rightarrow \quad
x = r\cos\alpha
\]

\begin{center}
\begin{tikzpicture}[scale=0.72]
\coordinate (O) at (0,0);
\coordinate (X) at (4,0);
\coordinate (P) at (4,2.5);

\draw[->] (-0.5,0) -- (5,0) node[right] {$x$};
\draw[->] (0,-0.5) -- (0,3.1) node[above] {$y$};

\draw[thick,blue] (O)--(P) node[midway,above left] {$r$};
\draw[very thick,orange] (O)--(X) node[midway,below] {$x=r\cos\alpha$};
\draw[dashed] (X)--(P);

\draw[->,red,thick] (0.9,0) arc (0:32:0.9);
\node[red] at (1.2,0.35) {$\alpha$};

\filldraw[blue] (P) circle (2pt);
\node[above right] at (P) {$P$};

\node[orange] at (2,-0.55) {Proyección horizontal};
\end{tikzpicture}
\end{center}

\textbf{Idea:} el coseno indica cuánto del radio \(r\) se proyecta sobre el eje horizontal.
\end{frame}

%---------------------------------
\begin{frame}{Origen del seno}
\small
El seno se define como la razón entre el cateto opuesto y la hipotenusa:

\[
\sin\alpha = \frac{y}{r}
\quad \Rightarrow \quad
y = r\sin\alpha
\]

\begin{center}
\begin{tikzpicture}[scale=0.72]
\coordinate (O) at (0,0);
\coordinate (X) at (4,0);
\coordinate (P) at (4,2.5);

\draw[->] (-0.5,0) -- (5,0) node[right] {$x$};
\draw[->] (0,-0.5) -- (0,3.1) node[above] {$y$};

\draw[thick,blue] (O)--(P) node[midway,above left] {$r$};
\draw[dashed] (O)--(X);
\draw[very thick,green!60!black] (X)--(P) node[midway,right] {$y=r\sin\alpha$};

\draw[->,red,thick] (0.9,0) arc (0:32:0.9);
\node[red] at (1.2,0.35) {$\alpha$};

\filldraw[blue] (P) circle (2pt);
\node[above right] at (P) {$P$};

\node[green!50!black] at (4.85,1.25) {Proyección vertical};
\end{tikzpicture}
\end{center}

\textbf{Idea:} el seno indica cuánto del radio \(r\) se proyecta sobre el eje vertical.
\end{frame}

%---------------------------------
\begin{frame}{Relación completa}
\small
Como el punto \(P\) se obtiene mediante sus dos proyecciones:

\[
x = r\cos\alpha
\qquad
y = r\sin\alpha
\]

entonces:

\[
P = (x,y) = (r\cos\alpha,\; r\sin\alpha)
\]

\begin{center}
\begin{tikzpicture}[scale=0.68]
\coordinate (O) at (0,0);
\coordinate (X) at (4,0);
\coordinate (P) at (4,2.5);

\draw[->] (-0.5,0) -- (5,0) node[right] {$x$};
\draw[->] (0,-0.5) -- (0,3.1) node[above] {$y$};

\draw[thick,blue] (O)--(P) node[midway,above left] {$r$};
\draw[very thick,orange] (O)--(X) node[midway,below] {$r\cos\alpha$};
\draw[very thick,green!60!black] (X)--(P) node[midway,right] {$r\sin\alpha$};

\draw[->,red,thick] (0.9,0) arc (0:32:0.9);
\node[red] at (1.2,0.35) {$\alpha$};

\filldraw[blue] (P) circle (2pt);
\node[above right] at (P) {$P=(r\cos\alpha,r\sin\alpha)$};
\end{tikzpicture}
\end{center}

Esta ecuación conecta geometría, trigonometría y coordenadas cartesianas.
\end{frame}

%---------------------------------
\begin{frame}{¿Qué significa rotar?}
Rotar un punto significa:

\[
\alpha \rightarrow \alpha + \theta
\]

\begin{block}{Idea clave}
No cambia la distancia al origen; solo cambia el ángulo.
\end{block}

Por eso una rotación es un movimiento sobre una circunferencia.
\end{frame}

%---------------------------------
\begin{frame}{Visualización de la rotación}
\begin{center}
\begin{tikzpicture}[scale=0.8]
\draw[->] (-0.5,0) -- (5,0) node[right] {$x$};
\draw[->] (0,-0.5) -- (0,4) node[above] {$y$};

\coordinate (P) at (4,1.5);
\coordinate (Q) at (2.2,3.5);

\draw[thick,blue] (0,0)--(P) node[midway,below] {$r$};
\draw[thick,purple] (0,0)--(Q) node[midway,left] {$r$};

\filldraw[blue] (P) circle (2pt);
\filldraw[purple] (Q) circle (2pt);

\node[right] at (P) {$P$};
\node[above] at (Q) {$P'$};

\draw[->,red,thick] (1,0) arc (0:20:1);
\draw[->,purple,thick] (1.3,0.5) arc (20:60:1.5);

\node[red] at (1.2,0.3) {$\alpha$};
\node[purple] at (1.6,1) {$\theta$};
\end{tikzpicture}
\end{center}

El punto \(P\) gira hasta \(P'\), pero conserva la misma distancia \(r\).
\end{frame}

%---------------------------------
\begin{frame}{Aplicación trigonométrica}
Después de rotar:

\[
P'=(r\cos(\alpha+\theta),\; r\sin(\alpha+\theta))
\]

Por tanto:

\[
x' = r\cos(\alpha+\theta)
\]

\[
y' = r\sin(\alpha+\theta)
\]

\medskip

Ahora necesitamos expresar \(x'\) y \(y'\) usando \(x\) y \(y\).
\end{frame}

%---------------------------------
\begin{frame}{Identidades trigonométricas: punto de partida}
Después de rotar, el punto queda descrito por:

\[
x'=r\cos(\alpha+\theta)
\]
\[
y'=r\sin(\alpha+\theta)
\]

Para separar el ángulo compuesto \(\alpha+\theta\), usamos:

\[
\cos(\alpha+\theta)=
\cos\alpha\cos\theta-\sin\alpha\sin\theta
\]

\[
\sin(\alpha+\theta)=
\sin\alpha\cos\theta+\cos\alpha\sin\theta
\]
\end{frame}

%---------------------------------
\begin{frame}{Derivación de \(x'\): primer paso}
Partimos de la coordenada horizontal rotada:

\[
x'=r\cos(\alpha+\theta)
\]

Sustituimos la identidad del coseno de una suma:

\[
x'=r(\cos\alpha\cos\theta-\sin\alpha\sin\theta)
\]

Distribuimos el factor \(r\):

\[
x'=r\cos\alpha\cos\theta-r\sin\alpha\sin\theta
\]
\end{frame}

%---------------------------------
\begin{frame}{Derivación de \(x'\): sustitución cartesiana}
Recordemos la relación polar--cartesiana:

\[
x=r\cos\alpha
\]
\[
y=r\sin\alpha
\]

En la expresión:

\[
x'=r\cos\alpha\cos\theta-r\sin\alpha\sin\theta
\]

reemplazamos \(r\cos\alpha\) por \(x\), y \(r\sin\alpha\) por \(y\):

\[
x'=x\cos\theta-y\sin\theta
\]
\end{frame}

%---------------------------------
\begin{frame}{Derivación de \(y'\): primer paso}
Partimos de la coordenada vertical rotada:

\[
y'=r\sin(\alpha+\theta)
\]

Sustituimos la identidad del seno de una suma:

\[
y'=r(\sin\alpha\cos\theta+\cos\alpha\sin\theta)
\]

Distribuimos el factor \(r\):

\[
y'=r\sin\alpha\cos\theta+r\cos\alpha\sin\theta
\]
\end{frame}

%---------------------------------
\begin{frame}{Derivación de \(y'\): sustitución cartesiana}
Recordemos nuevamente:

\[
x=r\cos\alpha
\]
\[
y=r\sin\alpha
\]

En la expresión:

\[
y'=r\sin\alpha\cos\theta+r\cos\alpha\sin\theta
\]

reemplazamos \(r\sin\alpha\) por \(y\), y \(r\cos\alpha\) por \(x\):

\[
y'=y\cos\theta+x\sin\theta
\]

Equivalentemente:

\[
y'=x\sin\theta+y\cos\theta
\]
\end{frame}

%---------------------------------
\begin{frame}{Resultado final de la rotación}
Las dos coordenadas después de rotar son:

\[
x'=x\cos\theta-y\sin\theta
\]

\[
y'=x\sin\theta+y\cos\theta
\]

\begin{block}{Lectura geométrica}
La nueva coordenada horizontal y la nueva coordenada vertical se obtienen combinando las proyecciones de \(x\) y \(y\) según el ángulo de rotación \(\theta\).
\end{block}
\end{frame}


%---------------------------------
\begin{frame}{Forma matricial}
\[
\begin{bmatrix}
x' \\
y'
\end{bmatrix}
=
\begin{bmatrix}
\cos\theta & -\sin\theta \\
\sin\theta & \cos\theta
\end{bmatrix}
\begin{bmatrix}
x \\
y
\end{bmatrix}
\]

\medskip

La matriz encapsula toda la transformación de rotación.
\end{frame}


%---------------------------------
\begin{frame}{Interpretación}
La rotación:

\begin{itemize}
\item Mezcla las coordenadas \(x\) y \(y\).
\item Usa seno y coseno como coeficientes de mezcla.
\item Conserva la longitud del vector.
\item Cambia la orientación del punto.
\end{itemize}

\medskip

Por eso la rotación es una transformación rígida.
\end{frame}

%---------------------------------
\begin{frame}{Conclusión}
\begin{itemize}
\item Seno y coseno describen proyecciones.
\item La rotación es movimiento sobre un círculo.
\item La fórmula surge de la geometría.
\item La matriz organiza las ecuaciones de rotación.
\end{itemize}

\medskip

\textbf{Idea final:} la rotación no se memoriza; se comprende geométricamente.
\end{frame}

\end{document}
